\documentclass{article}

\usepackage{amsmath, makecell}

\title{Ch. 25 HW \#1 Solutions}
\author{Principles of Macroeconomics}
\date{\today}

\begin{document}

\maketitle

\section*{Problems \& Applications}

\begin{enumerate}

\setcounter{enumi}{1}

\item

	\begin{enumerate}
	
	\item \phantom{x}
		
		\begin{center}
		\begin{tabular}{| c | c | c | c |}
		\hline
		Year & $P_{Cauliflower}$ & $P_{Broccoli}$ & $P_{Carrots}$ \\
		\hline
		2023 & $200 / 100 = \$2$ & $75 / 50 = \$1.50$ & $50 / 500 = \$0.10$ \\
		\hline
		2024 & $225 / 75 = 3$ & $120 / 80 = 1.5$ & $100 / 500 = 0.2$ \\
		\hline
		\end{tabular}
		\end{center}
		
	\item I'm going to assume that the typical consumption basket is 4 heads of Cauliflower, 2 bunches of broccoli, and 20 carrots. Your basket might be different. We aren't told what to assume, so you can assume whatever you want.
	
		\begin{tabular}{| c | c | c |}
		\hline
		Year & \makecell{Cost of\\Basket} & CPI \\
		\hline
		2023 & $4(2) + 2(1.5) + 20(0.1) = 13$ & 100 \\
		\hline
		2024 & $4(3) + 2(1.5) + 20(0.2) = 19$ & $19 / 13 * 100 = 146$ \\
		\hline
		\end{tabular}
		
	\item  $\frac{146 - 100}{100} \cdot 100 = 46\%$
		
	\end{enumerate}
	
\item
	
	\begin{enumerate}
		
	\item Tennis Balls: 0 \\
		Golf Balls: $\frac{6 - 4}{4} \cdot 100 = 50\%$ \\
		Gatorade: $\frac{2 - 1}{1} \cdot 100 = 100\%$
		
	\item We're not told how to fix the basket, so it's up to you. I'm going to assume the typical person buys one tennis ball, one golf ball, and two bottles of gatorade.
		
		\begin{tabular}{| c | c | c | c |}
		\hline
		Year & \makecell{Cost of\\Basket} & CPI & Inflation \\
		\hline
		2023 & $1(2) + 1(4) + 2(1) = 8$ & 100 & -- \\
		\hline
		2024 & $1(2) + 1(6) + 2(2) = 12$ & $12 / 8 * 100 = 150$ & $\frac{150 - 100}{100} \cdot 100 = 50\%$ \\
		\hline
		\end{tabular}
			
	\item It should. There's been a quality improvement in Gatorade bottles, so the price increase is partially due to quality improvement and not just an increase in the cost of living. As a result, our inflation rate should be lower than the one we calculated in (b).
		
	\item It should. The introduction of new products increases the value of a dollar. That increased value offsets some of the increase in the cost of living, so our inflation rate should be lower than the one we calculated in (b).
		
	\end{enumerate}
		
\item You don't need me for this one.
	
\item 

	\begin{enumerate}
	
	\item \phantom{x}
		
		\begin{center}
		\begin{tabular}{| c | c | c |}
		\hline
		Year & \makecell{Cost of\\Basket} & CPI \\
		\hline
		2023 & $5(40) + 3(50) = 350$ & 100\\
		\hline
		2024 & $5(60) + 3(60) = 480$ & $\frac{480}{350} \cdot 100 = 137$ \\
		\hline
		\end{tabular}
		\end{center}
		
		Inflation rate $= 37\%$
		
	\item \phantom{x}
	
		\begin{center}
		\begin{tabular}{| c | c | c | c |}
		\hline
		Year & $GDP_{Nom}$ & $GDP_{Real}$ & $GDP_{Def}$ \\
		\hline
		2023 & $10(40) + 6(50) = 700$ & 700 & 100 \\
		\hline
		2024 & $12(60) + 10(60) = 1,320$ & $12(40) + 10(50) = 980$ & $\frac{1,320}{980} \cdot 100 = 135$ \\
		\hline
		\end{tabular}
		\end{center}
		
		Inflation rate $= 35\%$
		
	\item It is not. The CPI and the GDP deflator weight prices differently. The CPI weights them based on a fixed basket that approximates what consumers typically buy. The GDP deflator weights them based on how much of the good was produced.
	
	\end{enumerate}
	
\item

	\begin{enumerate}
	
	\item The introduction of cell phones is the introduction of a new good which increases the value of a dollar and decreases the cost of living. Since the CPI uses a fixed basket, it doesn't account for new goods and overstates the cost of living.
	
	\item The introduction of airbags in cars is a quality improvement in cars. If the addition of airbags is unmeasured and increases the cost of cars, it will increase the CPI even though it shouldn't because the increase in price reflects the improved safety of cars, not a devaluation of the dollar.
	
	\item Increased personal computer purchases in response to a decline in their price leads to substitution bias in the CPI. The change in consumer behavior creates a discrepancy between the fixed basket underlying the CPI and the basket that consumers are actually buying. As a result, the CPI won't accurately reflect the cost of living. 
	
	\item More scoops of raisins in each package of Raisin Bran is a quality improvement. See explanation above.
	
	\item Greater use of fuel-efficient cars after gas prices increase leads to substitution bias in the CPI. See explanation above. 
	
	\end{enumerate}
			
\end{enumerate}

\end{document}